\documentclass[a4paper,12pt]{article}
\usepackage{graphicx}
\usepackage[a4paper, top=0cm, bottom=2cm, left=2cm, right=2cm]{geometry}

\usepackage{xcolor}
\usepackage{hyperref}
\usepackage{eso-pic}

\usepackage[french]{babel}
\usepackage[T1]{fontenc}
\usepackage{amsmath, amsfonts, amssymb, mathrsfs}

\newcounter{correction}

\newcommand{\question}{%
  \refstepcounter{correction}%
  \textbf{\textcolor{black}{Question \thecorrection\ (1 point) :}} \ignorespaces
}

% Filigrane
\AddToShipoutPicture{
    \AtTextCenter{%
        \makebox[0pt]{\rotatebox{80}{\textcolor[gray]{0.7}{\fontsize{5cm}{5cm}\selectfont PGB}}}
    }
}

\begin{document}

\hrule
\begin{center}
    \begin{tabular}{@{} p{5cm} p{5cm} p{5cm} @{}} 
        Lycée Dindéfélo & \quad\quad Test 1 & 00 Mois 2026 \\
    \end{tabular}
    \\[-0.01cm]
    \hrule
\end{center}

\begin{center}
    \textbf{\Large Calcul Numérique} \\[0.2cm]
    \textbf{\large Professeur : M. BA} \\[0.2cm]
    \textbf{Classe : Terminale S2} \\[0.2cm]
    \textbf{\small Durée : 30 minutes}\\[0.2cm]
    \textbf{\small Note : \quad /5}
\end{center}

\textbf{\small Nom de l'élève :} \underline{\hspace{8cm}} \\[0.7cm]

\section*{Exercice 1 : Ordre de priorité (Niveau 6ème)}
\textit{Calculer les expressions suivantes en soulignant à chaque étape le calcul prioritaire.}

\begin{enumerate}
    \item $A = 25 - 5 + 3$
    \item $B = 24 \div 4 \times 2$
    \item $C = 12 + 3 \times 5$
    \item $D = 40 - (8 + 3 \times 2)$
\end{enumerate}

\vspace{1cm}

\section*{Exercice 2 : Parenthèses et crochets (Niveau 5ème)}
\textit{Calculer les expressions suivantes en respectant les priorités (parenthèses, puis crochets).}

\begin{enumerate}
    \item $E = 15 - [7 - (3 + 1)]$
    \item $F = 3 \times [12 - (5 - 2)]$
    \item $G = \dfrac{18 + 6}{5 - 2}$
\end{enumerate}

\vspace{1cm}

\section*{Exercice 3 : Nombres relatifs et règles de signes (Niveau 4ème)}
\textit{Effectuer les calculs suivants en appliquant les règles de signes.}

\begin{enumerate}
    \item $H = (-8) + (+5) - (-3) + (-12)$
    \item $I = (-4) \times (+6)$
    \item $J = (-5) \times (-3) \times (-2)$
    \item $K = \dfrac{-24}{-4}$
    \item $L = (+15) - [(-7) - (+3) \times (-2)]$
\end{enumerate}

\vspace{1cm}

\section*{Exercice 4 : Bilan et complexité (Niveau 4ème / Approfondissement)}
\textit{Calculer l'expression suivante en détaillant chaque étape :}

$$M = -10 + \frac{15 - 3 \times 5}{-2} - [(-4) \times (-2)]$$

\end{document}